% Synthetic test fixture for Alexandria's PDF extraction tests.
%
% Built with default Computer Modern so the compiled PDF embeds the cmmi/cmsy/cmex
% math fonts that the math-font detector (_has_math_fonts) keys on. It deliberately
% uses no hyperref/\pdfinfo title, so the PDF carries no /Title metadata (mirroring
% typical pdfTeX output, which the extraction test asserts on).
%
% Regenerate reproducibly (fixed 2025 /CreationDate) with:
%   SOURCE_DATE_EPOCH=$(date -u -d 2025-08-08 +%s) FORCE_SOURCE_DATE=1 \
%     pdflatex -interaction=nonstopmode paper.tex
%
% All content below is fictional and authored for this fixture.
\documentclass[11pt]{article}

\title{A Sample Paper on Vector Fields}
\author{Ada Fixtura \\ Alexandria Test Fixtures}
\date{2025}

\begin{document}
\maketitle

\begin{abstract}
This is a synthetic document used only to exercise Alexandria's PDF extraction
pipeline. It carries a few genuine equations so the compiled PDF embeds Computer
Modern math fonts, and enough prose to make chunking meaningful.
\end{abstract}

\section{Introduction}

We consider a toy magnetostatic field so the text contains real mathematics
rather than placeholder symbols. The magnetic field $\vec{B}$ produced by a
steady current follows the Biot--Savart law, and its axial component $B_z$ can
be written compactly in terms of a scalar potential $\Phi$. The Greek letters
and operators below guarantee a nonzero math-symbol density, and the equations
force the math fonts into the font resource dictionary.

\section{A toy field equation}

The field of a current loop of current $I$ is
\begin{equation}
\vec{B}(\vec{r}) = \frac{\mu_0}{4\pi} \oint \frac{I\, d\vec{l} \times \hat{r}}{r^2},
\end{equation}
and along the symmetry axis the component reduces to
\begin{equation}
B_z = \alpha\, \frac{\partial \Phi}{\partial z}, \qquad
\alpha = \frac{\mu_0 I}{2}.
\end{equation}
Here $\alpha$ is a constant, $\Phi$ is the scalar potential, and the summation
$\sum_{n=1}^{\infty} \frac{1}{n^2} = \frac{\pi^2}{6}$ appears purely to add a
few more symbols to the page.

\section{Discussion}

The point of this fixture is not physics but coverage. A real extractor must
survive a multi-page, \LaTeX-generated PDF: hyphenated line breaks, ligatures
such as \textit{efficient} and \textit{affordable}, and inline math mixed with
ordinary prose. None of the numbers here mean anything; they exist so that the
normalizer, the chunker, and the math-density heuristics all have something
representative to chew on.

We repeat a little more prose to push the document onto a second page, because
the extraction test treats this as a multi-page paper. Retrieval-augmented
generation, vector databases, and keyword search over an inverted index are all
mentioned here so the fixture reads like a plausible technical note rather than
lorem ipsum. The conclusion, such as it is, is that vector fields make excellent
excuses for embedding math fonts.

\end{document}
